\section{Limitations and Discussion}

\textbf{Code semantics are insufficient.} The controlled SECDED pair shows that exact-equivalent coding behavior does not uniquely determine implementation behavior. Register placement changes area, interconnect, latency, Fmax, and activity, creating a non-dominated trade rather than one universally better implementation.

\textbf{Stronger correction carries implementation cost.} The BCH point corrects all double-coordinate errors that SECDED only detects, but its evaluated realization is larger, more routing-intensive, and infeasible at the shared target. The comparison intentionally retains the extra six codeword bits as part of the realized cost, so it cannot attribute every difference solely to decoder algebra. Different BCH architectures could occupy different points~\cite{wong2011,lagendijk2026}.

\textbf{Missing data must remain missing.} Hsiao illustrates why qualification boundaries matter. Relaxing the synthesis-memory policy only after observing failure would create a different experiment; imputing PPA would be less defensible still. Explicit missingness narrows the conclusion without suggesting that Hsiao is physically inferior.

The empirical scope is limited to specific RTL implementations, SKY130HD, one TT 1.8 V/25~$^{\circ}$C corner, one common 10 ns target, and one deterministic seed/worker policy. Only four reliability identities, three routed designs, and two timing-feasible energy points are available. OpenSTA activity-based power is an estimate, not a silicon measurement. BCH target-clock energy is excluded, and Hsiao PPA is unavailable. We do not study process or seed distributions, multi-corner signoff, physical SRAM arrays, bit mapping, interleaving, scrubbing, or memory-macro cost. No FIT, SER, field weighting, or operational SDC/DUE probability is established. These limits bound generalization while leaving the controlled within-flow contrasts intact.

\subsection{Threats to Validity}

\textit{Internal validity.} A shared policy reduces configuration confounding but does not remove heuristic implementation effects. One seed and one worker provide a controlled point, not a distribution over place-and-route outcomes. The prospective no-retry rule prevents selectively polishing a poor result, while also leaving open whether a different equally fair policy would alter the ordering. Exact SECDED equivalence and the common transaction contract strengthen the microarchitecture contrast; the BCH comparison lacks such an exact-equivalent architectural pair.

\textit{Construct validity.} Final cell area, routed wirelength, post-route timing, and activity-based power operationalize implementation cost, but do not cover every system cost. The experiment excludes SRAM macro area, ECC storage organization, clock-domain crossings, control integration, and error-recovery latency. The equal-work trace supports a throughput-mode energy comparison, not energy under arbitrary workloads. Similarly, canonical-coordinate W3 enumeration measures decoder outcomes in a finite logical universe, not radiation susceptibility.

\textit{External validity.} A result from one open 130-nm library and one corner need not transfer quantitatively to advanced nodes, other cell libraries, memories, or decoder generators. The conclusion with the broadest support is therefore methodological: preserve RTL identity, guarantees, feasibility, and missingness when joining reliability with physical evidence. The numerical trade-offs remain bounded to the named artifacts.

\subsection{Implications for ECC Selection}

The evidence suggests a two-stage decision rather than a single ranking. First, establish non-negotiable requirements: correction guarantee, target clock, transaction latency, and acceptable evidence status. Second, compare surviving implementation identities on physical objectives. In this experiment, a strict W2-correction requirement admits BCH semantically but exposes a timing implementation gap; a strict 100 MHz requirement admits the two SECDED points but not BCH; and a three-cycle latency prohibition favors combinational SECDED despite the pipeline's energy and Fmax advantages. Reporting these conditional choices is more actionable than naming one ``best ECC.''

For designers, the missing Hsiao point is also informative. It identifies a representation boundary---an inferred table under a fixed synthesis-memory policy---that future work could address prospectively with an algorithmic encoder/decoder RTL. Such a future implementation would be a new identity and should be requalified, rather than used retroactively to fill the present table.

\subsection{Auditability and Failure-Closed Reporting}

Cross-layer studies are vulnerable to accidental evidence substitution: a proof from one RTL can be joined to PPA from another, a pre-route estimate can be displayed beside post-route measurements, or a power report can survive after its clock point becomes infeasible. Our artifact guards the joins with stable implementation identifiers and evidence status. A physical row is admitted only when its identity matches the qualified RTL and the common flow policy. The same rule propagates to derived quantities: Fmax retains measured WNS, and target-clock energy exists only when that target is feasible.

The presentation has three additional safeguards. First, guarantee statements and exhaustive W3 observations occupy different fields and are described with different language. Second, unavailable values remain textual status, not zeros, infinities, or coordinates in a plot. Third, full-precision evidence feeds generated macros, CSV tables, and plots, while manuscript rounding is display-only. The reported 36.702\%, 45.911\%, and 24.016\% SECDED changes, for example, are recomputed from the frozen raw values rather than copied from captions.

This discipline changes what a reviewer can falsify. The SECDED equivalence claim can be checked independently of physical cost; the BCH timing status can be recomputed from target period and WNS; the power comparison can be audited against operation count, trace duration, II, and feasibility; and every omitted point has a stated reason. None of these checks proves external validity, but they make the implementation-scoped claims inspectable and prevent an unsupported value from silently strengthening the result.

\subsection{What Would Change the Decision}

The present Pareto structure is contingent rather than immutable. A prospectively defined Hsiao RTL that passes the unchanged qualification and physical policy could add a new point without altering the current measurements. An alternative BCH architecture could meet 10 ns and thereby become eligible for equal-work target-clock energy, but its result could not be substituted for the evaluated implementation. Multi-seed or multi-corner characterization could reveal distributions around each point, while a physically grounded fault model could supply weights for the W3 masks. Each extension answers a distinct uncertainty and should preserve the same identity and feasibility rules.

These extensions also clarify which conclusions are already stable. No future point can make the two measured SECDED implementations code-inequivalent, erase their observed area/Fmax/energy trade-off under this environment, or turn the measured BCH WNS into a met constraint. Future evidence can broaden the design space and alter a selection, but should not rewrite the status of the frozen points.
